\section{From one spectrum to a closed surface}
\label{sec:dynamics}

The remaining step in the diagram induction concerns a closed facet.
Cutting a nonseparating circle reduces it to a facet with boundary, but
prescribes the spectrum of the gluing monodromy.  We first obtain density
with this spectrum fixed.  A Dehn twist then supplies the missing
directions.

Multiplication by a centralizer element is the algebraic form of
the generalized twist motions described by
Goldman~\cite[Section~4.6]{Goldman1986}.  We use the matrix relation
directly, then verify the needed torus closure and tangent directions
at an explicit representation.

\begin{lemma}\label{lem:closed-density}
Let $2\le n\le r$, $g\ge2$, and let $c\in F$ satisfy $c^n=1$.
Assume that integral solutions are dense for every normalized diagram
of ranks at most $n$ whose rank-$n$ facets have nonempty boundary.
Then $\mathcal R_{n,g}(c)(R)$ is Zariski dense in
$\mathcal R_{n,g}(c)$.
\end{lemma}

\begin{proof}
Write $V=\mathcal R_{n,g}(c)$ and let $D$ be the Zariski closure of its
integral points.  By Proposition~\ref{prop:surface-irreducibility}, $V$
is irreducible.

\smallskip\noindent\emph{Integral gluing at a fixed spectrum.}
Choose
\[
 \alpha_i=\zeta^{\,i-1}\qquad(1\le i\le n).
\]
These numbers are distinct and their pairwise differences are units of
$R$, by Lemma~\ref{lem:unit-spectrum}.  We first show that $D$ contains
the entire fiber
\[
 V_\alpha=\{(A_j,B_j)\in V:
            \det(XI-A_1)=\prod_{i=1}^n(X-\alpha_i)\}.
\]

Represent $V$ by a genus-$g$ facet with one puncture of scalar
monodromy $c^{-1}$, so that the surface relation is
$\prod_j[A_j,B_j]=cI$.
Cut this surface along the circle represented by $A_1$.  On each new
boundary, require a full flag with successive scalar monodromies
$\alpha_1,\ldots,\alpha_n$, using rank-one punctured disk caps.  Here
the two loops are read in the same original direction; their induced
boundary orientations are opposite, so one reverses the cap scalar
when necessary.  The new seams have $n>1$ rank-one children, so they
are proper; the resulting rank-$n$ facet has boundary.  The
hypothesis gives density of integral solutions on this cut diagram.

To reglue, choose an isomorphism between the two boundary systems.
Because the eigenvalues are distinct, the choices form a torsor for
$(\mathbb G_m)^n$, by matching corresponding eigenlines.  They vary
locally algebraically with the cut solution, so their projection to
the cut solution variety is open and surjective.
This description also holds over $R$.  Indeed the spectral projectors
\begin{equation}\label{eq:spectral-projectors}
 E_i(A)=\prod_{j\ne i}\frac{A-\alpha_jI}{\alpha_i-\alpha_j}
 \in M_n(R)
\end{equation}
split off the eigenmodules.  Write $F_\bullet$ for the prescribed
flag.  The projector $E_i$ is zero on $F_{i-1}$ and on $R^n/F_i$:
on each of their graded pieces its action is zero, so its induced
action is both nilpotent and idempotent.  It induces the identity on
$F_i/F_{i-1}$.  Hence its image is isomorphic to that graded line,
which is free because it is identified with a framed cap fiber.
Thus every integral cut
solution has gluing choices, and in these bases the integral choices
are $(R^\times)^n$, a dense subset of the gluing torus by
Lemma~\ref{lem:integral-parameters}.
The open-projection density principle therefore gives dense integral
glued solutions.  Conversely, every point of $V_\alpha$ admits these
ordered eigenflags and cap frames, and is recovered by gluing, up to
frame changes.  Frame changes over $R$ are dense by the same lemma.  Hence
$V_\alpha\subset D$.

\smallskip\noindent\emph{Torus saturation by Dehn twists.}
The automorphisms
\[
 \tau^{\pm1}(A_1,B_1,A_2,B_2,\ldots)
       =(A_1B_1^{\pm1},B_1,A_2,B_2,\ldots)
\]
preserve $V$, since $[A_1B_1,B_1]=[A_1,B_1]$, and preserve its integral
points.  Consequently they preserve $D$.  We show that this discrete
invariance enlarges $V_\alpha$ to an analytic open subset of $V$.

Choose distinct rational primes $\beta_1,\ldots,\beta_n$.  They are
multiplicatively independent.  Let $P$ be the cyclic permutation
matrix $Pe_i=e_{i+1}$, with indices modulo $n$.  The point $x_0\in V$
given by
\begin{equation}\label{eq:twist-test-point}
 \begin{aligned}
 A_1&=\operatorname{diag}(\alpha_1,\ldots,\alpha_n),&
 B_1&=\operatorname{diag}(\beta_1,\ldots,\beta_n),\\
 A_2&=P,&
 B_2&=\operatorname{diag}(1,c^{-1},\ldots,c^{1-n})
 \end{aligned}
\end{equation}
and identity matrices on the other handles belongs to $V_\alpha$.
Indeed $PB_2P^{-1}=cB_2$, including the wraparound entry because
$c^n=1$.  The numbers $\beta_i$ are used only to choose this complex
point; they need not be units of $R$.

The relation map $\GL_n^{2g}\to\SL_n$ is submersive at $x_0$.
Left infinitesimal variations of $B_1$ give
$(\operatorname{Ad}_{A_1}-1)X$, spanning all off-diagonal matrices.
Diagonal variations of $B_2$ give
$(\operatorname{Ad}_P-1)H$, spanning the traceless diagonal matrices.
These span $\mathfrak{sl}_n$, so $V$ is smooth at $x_0$.
The map $\sigma$ recording the characteristic coefficients of $A_1$
is also submersive there: varying $A_1$ diagonally preserves the
relation, and the elementary-symmetric-functions map has invertible
differential at distinct eigenvalues.  Thus $V_\alpha$ is smooth near
$x_0$.

Take a small connected analytic neighborhood $U$ of $x_0$ in
$V_\alpha$.  The eigenvalues $b_i(x)$ of $B_1(x)$ and their spectral
projectors $P_i(x)$ can be labeled holomorphically on $U$, since the
spectrum at $x_0$ is simple.  For each nonzero $m\in\Z^n$, the
holomorphic function
\[
             \prod_i b_i(x)^{m_i}-1
\]
is nonzero at $x_0$.  Its zero set is therefore closed with empty
interior.  The Baire theorem implies that the points $x\in U$ with
multiplicatively independent $b_i(x)$ form a dense subset $U_0$.
For such an $x$, the powers of $B_1(x)$ are Zariski dense in its
centralizer torus: otherwise a nontrivial character of that torus
would equal one on the cyclic subgroup, giving a multiplicative relation.
Since $x\in D$ and $\tau^k(x)\in D$ for every $k\in\Z$, closedness
then gives every tuple obtained by replacing $A_1(x)$ with $A_1(x)T$,
for $T$ in this centralizer.

\smallskip\noindent\emph{A submersive family of twists.}
To extend this conclusion over all of $U$, write the centralizers in
the holomorphic form
\[
       T(x,t)=\sum_{i=1}^n t_iP_i(x),
       \qquad t\in(\C^\times)^n.
\]
Replacing $A_1(x)$ with $A_1(x)T(x,t)$ defines a holomorphic map
$\Psi:U\times(\C^\times)^n\to V$; it preserves the relation because
$T(x,t)$ commutes with $B_1(x)$.  Its values on
$U_0\times(\C^\times)^n$ lie in $D$.  Since $D$ is analytically
closed, all values of $\Psi$ lie in $D$.

At $(x_0,1)$ the directions from $U$ give $\ker(d\sigma)$.
The torus directions independently multiply the diagonal entries of
$A_1$, so their images under $d\sigma$ span its target.  Hence
$d\Psi$ is surjective.  Its image contains an analytic open subset of
$V$.  Such a subset is Zariski dense in an irreducible complex
variety, proving $D=V$.
\end{proof}

\subsection{Completion of the rank induction}

\begin{proof}[Proof of Theorem~\ref{thm:diagram-density}]
We induct on the largest facet rank $n$. The rank-zero diagram has
only empty data. Normalize zero pieces and identity seams as in
Section~\ref{subsec:identity-normalization}; solutions of the original
diagram are recovered from solutions of the normalized diagram by
integral restriction and frame changes. At rank one no proper seam
can remain, since it would have at least two positive-rank children.
Thus the normalized diagram consists of independent closed scalar
local systems. Each solution variety is empty or a torus, so the
rank-one case follows from Lemma~\ref{lem:integral-parameters}.

Now let $n\geq2$ and assume the theorem for all diagrams of ranks
less than $n$. Proposition~\ref{prop:nonclosed-reduction} proves
density for the part of the normalized diagram having no closed
rank-$n$ facet. It allows all lower-rank facets, closed or not.
Every remaining closed rank-$n$ facet is an independent factor
$\mathcal R_{n,g}(c)$, where $c$ is the inverse product of its scalar
puncture monodromies. If $c^n\ne1$, that factor is empty. Otherwise
the genus-zero factor is a point or empty,
Lemma~\ref{lem:torus-density} gives density for genus one, and
Lemma~\ref{lem:closed-density} gives it for higher genus. The
hypothesis of the last lemma is exactly the normalized nonclosed
step already established at this rank.

The product of these dense sets is dense in the normalized solution
variety. Restoring the original diagram and its frames preserves
density by Lemma~\ref{lem:integral-parameters}. This completes the
induction over the same fixed ring $R$.
\end{proof}
